%!TEX root = ED_Grafo.tex

\section{Conceito}
\begin{frame}[label=notas]
	Grafo: uma representação de relações (arestas) entre objetos (vértices).
	\vfill\pause
	Um grafo $G=(V,A)$ é definido por:
	\begin{description}
		\item[$V$] conjunto de vértices (nós).
		\item[$A$] conjunto de arestas (arcos), que são pares [ordenados] de 
		vértices de $V$.
	\end{description}
	\vfill\pause
	Exemplo:
	\begin{itemize}
		\item[V] = \{A,B,C,D,E,F\}
		\item<4->[A] = \{(A,B), (B,C), (B,D), (D,A), (D,C), (E,F)\}
	\end{itemize}
	\vfill\pause
	\begin{center}
		\begin{tikzpicture}
			\drawABCDnodes%
			\drawEFnodes%
			\pause%
			\drawABCDedges%
			\drawEFedges%
		\end{tikzpicture}
	\end{center}
\end{frame}

\begin{frame}
	\framesubtitle{Aplicações}
	
	\pause
	\begin{itemize}
		\item[$G=$] (Cidades, Estradas)
		\item[$G=$] (Dispositivos Eletrônicos, Conexões)
		\item[$G=$] (Tarefas, Dependências)
		\item[$G=$] (Pessoas, Parentesco)
		\item[$G=$] (etc., etc.)
	\end{itemize}
\end{frame}